\documentclass[../../../include/open-logic-section]{subfiles}

\begin{document}

\olfileid[hi]{sth}{choice}{banach}
\olsection{बानाख--टार्स्की विरोधाभास}

जैसे बूलोस ने प्रतिस्थापन का औचित्य देने का प्रयास किया था, वैसे ही हम भी
\emph{बाह्य} विचारों का सहारा लेकर चयन का औचित्य देने का प्रयास कर सकते हैं
(\olref[replacement][extrinsic]{sec} देखें)। आखिर चयन को अपनाने के अनेक
वांछनीय परिणाम हैं: प्रत्येक कार्डिनल की तुलना करने की क्षमता; प्रत्येक
समुच्चय को सुव्यवस्थित करने की क्षमता; कार्डिनलों को एक विशिष्ट प्रकार का
क्रमसूचक मानने की क्षमता; इत्यादि।

किंतु कभी-कभी दावा किया जाता है कि चयन के \emph{अवांछनीय} परिणाम हैं। इसका
मुख्य कारण \cite{BanachTarski1924} का एक परिणाम है।

\begin{thm}[बानाख--टार्स्की विरोधाभास ($\ZFC$ में)]
किसी भी ठोस गोले को परिमिततः अनेक टुकड़ों में विघटित किया जा सकता है, जिन्हें
(घूर्णन और स्थानांतरण द्वारा) फिर जोड़कर उस गोले की दो प्रतियाँ बनाई जा सकती
हैं।
\end{thm}
\noindent 
पहली दृष्टि में यह कुछ विस्मयकारी है। स्पष्टतः दोनों गोलों का आयतन मूल गोले
के आयतन का \emph{दोगुना} है। किंतु दृढ़ गतियाँ---घूर्णन और स्थानांतरण---आयतन
नहीं बदलतीं। इसलिए ऐसा लगता है जैसे बानाख--टार्स्की हमें जादू से नया पदार्थ
अस्तित्व में लाने देता है।

बात और भी बिगड़ती है।\footnote{\citet[Theorem 3.12]{Wagon2016} देखें।}
समान तर्क दिखाता है कि एक मटर को परिमिततः अनेक टुकड़ों में काटा जा सकता है,
जिन्हें फिर (घूर्णन और स्थानांतरण द्वारा) जोड़कर बिग बेन के आकार और आकृति
वाली वस्तु बनाई जा सकती है।

किंतु इनमें से कुछ भी अकेले $\ZF$ में सिद्ध नहीं होता।\footnote{यद्यपि
बानाख--टार्स्की को चयन से पूर्णतः दुर्बल सिद्धांतों से सिद्ध किया जा सकता
है; \citet[303]{Wagon2016} देखें।} अतः हमारे सामने निर्णय है: चयन को
अस्वीकार करें या ``विरोधाभास'' के साथ जीना सीखें।

हम सुझाएँगे कि हमें इस ``विरोधाभास'' के साथ जीना सीखना चाहिए। वास्तव में
हमें नहीं लगता कि यह कोई विशेष विरोधाभास है। खासकर हमें नहीं दिखता कि यह
निम्न परिणामों में से किसी से अधिक या कम विरोधाभासी क्यों है:
\footnote{\citet[276--7]{Potter2004}, \citet[16]{Weston2003} और
\citet[31, 308--9]{Wagon2016} अन्य उदाहरणों से समान बातें कहते हैं।}
\begin{enumerate}
	\item अंतराल $(0,1)$ में उतने ही बिंदु हैं जितने $\Real$ में।
	\\\emph{प्रमाण}: $\tan(\pi(r-\nicefrac{1}{2})))$ पर विचार करें।
	\item किसी रेखा में उतने ही बिंदु हैं जितने किसी वर्ग में।
	\\\olref[his][set][pathology]{sec} और \olref[his][set][cantorplane]{sec} देखें।
	\item स्थान-पूरक वक्र हैं।
	\\\olref[his][set][pathology]{sec} और \olref[his][set][hilbertcurve]{sec} देखें।
\end{enumerate}
इन तीनों परिणामों में किसी को चयन की आवश्यकता नहीं है। वास्तव में अब हम
उन्हें गणित के आश्चर्यजनक, सुंदर अंश मात्र मानते हैं। शायद हमें
बानाख--टार्स्की विरोधाभास के प्रति भी यही दृष्टिकोण अपनाना चाहिए।

निस्संदेह यहाँ एक तकनीकी प्रेक्षण आवश्यक है; किंतु उसके लिए केवल शांत भाव
से सोचना चाहिए। दृढ़ गतियाँ आयतन सुरक्षित रखती हैं। फलतः जिन पाँच
\footnote{हमने विरोधाभास को ``परिमिततः अनेक टुकड़ों'' के पदों में कहा था।
वास्तव में \citet{Robinson1947} ने सिद्ध किया कि विघटन \emph{पाँच} टुकड़ों
से (पर उससे कम से नहीं) किया जा सकता है। प्रमाण के लिए
\citet[pp.~66--7]{Wagon2016} देखें।} टुकड़ों में गोला विघटित होता है, वे सभी
\emph{मापनीय} नहीं हो सकते। मोटे तौर पर इन अलग-अलग टुकड़ों को आयतन देना
अर्थहीन है। आपको इन्हें बिंदुओं के अचित्रणीय ``अनंत बिखराव'' के रूप में
सोचना चाहिए। संभव है कि ऐसे ``अनंततः बिखरे'' समुच्चयों की कल्पना ``विचित्र''
लगे। किंतु उनका अस्तित्व \stagesacc{} में निहित इस निर्देश से निष्पन्न होता
प्रतीत होता है कि पहले उपलब्ध समुच्चयों के \emph{सभी संभव} संग्रह बनाए जाएँ।

यदि इनमें से कुछ भी आश्वस्त न करे, तो बानाख--टार्स्की विरोधाभास अपनाने के
पक्ष में यह अंतिम (बाह्य) तर्क है। इससे तत्काल सर्वकालीन श्रेष्ठ गणितीय
चुटकुला मिलता है:
\begin{enumerate}
	\item[] \emph{प्रश्न}। ``Banach--Tarski'' का वर्णविपर्यय क्या है?
	\item[] \emph{उत्तर}। ``Banach--Tarski Banach--Tarski''।
\end{enumerate}

\end{document}
